\section{Simulation}
\label{sec:simulation}
In order to analyze the effectiveness of a time analysis attack in a
realistic scenario, we decide to consider a simulation environment.
We wanted to model the scenario in a way that most of the
characteristics of the Tor network architecture would have been
considered.

The \textbf{shadow simulator}\cite{shadow} seemed to be a good candidate for our purpose. In
the next paragraphs we will introduce shadow and we will illustrate our
work based on this simulator.
 
\subsection{Shadow}
Most of the simulators used for networks experimentation do not provide
the use of real external applications, which their
behaviour are often simulated as well. 
This approach can be reasonable if the analysis is focused on the 
network layers below the application layer. 

In our case we are interested in studying the characteristic of the Tor
network which is based on the network application layer.
Thus the Tor application is one of the core part of the Tor architecture as it implements the
onion routing protocol that let the Tor nodes communicate each others.

It is evident that the Tor application has a significant role in
the Tor architecture, thus we chose the shadow simulator as it permits
the use of the real Tor applications in a simulated environment.
Essentially it allows the execution of a set of local applications that
communicates in a simulated computer network. Moreover we could
avoid the oversimplification of the system that could have been occurred
with a custom implementation.

Applications are executed inside the shadow simulator through 
dynamic libraries called plug-ins. 
These plug-ins are interfaces used by shadow to trace a
selective set of system functions and re-route them to the simulator instead
of letting them proceed directly to the kernel. 
In this way, the simulator is transparent to the application
that may function as like it was running in a standard UNIX environment. 

The plug-in that allows the execution of the Tor applications is
scallion. The latter provides also some useful tools for the virtual
Tor network topology generation.

The virtual network is modeled with a XML configuration
file, called blueprint, used by shadow to understand the network
structure and some network properties such as link latency, jitter and
packet loss rate. The blueprint also tells Shadow what software each
virtual node should run at its creation. This is specified with a
plug-ins list for each virtual node entry of the XML file.

We implemented three shadow plug-ins and their relative applications 
 able to work alongside scallion in
order to experiment a time analysis attack scenario. Our plug-ins and
the attack scenario will be illustrated in the next paragraphs.

%TODO: find some pictures
\subsection{Simulation Scenario}
%TODO: un bello schema con i nodi e i relativi plug-in

\subsection{Autosys plug-in}
To analyze the tor network traffic we need, at least, some information about
the incoming and the outgoing connections.\\
To achieve this result we have more than a solution: in the first one we can place 
in the last network element that is under the control of an autonomous system\footnote{
That the reason why the plug-in is called "Autosys".} a connection
listener which traces a whole subset of
network nodes;
otherwise we can place the listener in a, malware-like, invisible proxy on the target
machine with the aim to trace every outgoing connection. 
In both solutions we need to place some correspondent exit tracers on
the target end points.
This proxies can be implemented as a simple raw-socket based sniffer. 
Unfortunately Shadow doesn't support raw-socket
emulation\footnote{If you launch a simulation with a raw socket based program the simulator
inform you about the impossibility of the task.},
so we implemented the plug-in as a simple event based TCP
proxy. In our attack scenario an instance of the autosys plug-in is placed 
between the client application and the client Tor daemon, also a second
 instance is placed between the server Tor daemon and the server application.
\\
\begin{algorithm}
\caption{Autosys main loop}
\begin{algorithmic}[1]
\Function{Autosys}{$sa$, $sb$, $analyzer$}
	\While{$True$}
		\If {Some data can be read from $sa$}
			\State $indata \gets $ readall($sa$)
			\State writeall($sb$, $indata$)
			\If {$indata$ was a connection initiator}
				\State send($<type(sa)$;$hostname(sa)$;$gettod()>$, analyzer)
			\EndIf
		\EndIf
		\If {Some data can be read from $sb$}
			\State $indata \gets $ readall($sb$)
			\State writeall($sa$, $indata$)
			\If {$indata$ was a connection initiator}
				\State send($<type(sa)$;$hostname(sa)$;$gettod()>$, analyzer)
			\EndIf
		\EndIf
	\EndWhile
\EndFunction
\end{algorithmic}
\label{alg:auto_main}
\end{algorithm}

As we can see in the algorithm \ref{alg:auto_main}, where $sa$ and $sb$ are the two end sockets and
$analyzer$ is the socket descriptor of the analysis server (the logger), when some connection
initiator packets flow trough the proxy the
current machine time of day is sent to the analysis server
% DA CONTROLLARE FOOTNOTE
\footnote{The simulator have a single clock, so we don't have any problem with the
clock synchronization.}.

\begin{figure}[H]
\centering
\begin{tikzpicture}[scale=2]
\node (app) [rectangle, draw, thick, rounded corners=3pt, text width=100pt, fill=blue!50]{Application};
\node (autosys) [rectangle, draw, thick, rounded corners=3pt, text width=100pt, fill=red!50, below = 0.5cm of app]{Autosys};
\node (tor) [rectangle, draw, thick, rounded corners=3pt, text width=100pt, below = 0.5cm of autosys]{\begin{center}TOR daemon\\\includegraphics[scale=0.3]{tor.jpg}\end{center}};

\node (internet) [cloud, draw, thick, fill=red!20!blue!20!green!20, right = 2 cm of tor] {Net};

\node (analysis) [rectangle, draw, thick, rounded corners=3pt, text width = 100pt, below= 1cm of internet, fill=yellow!50] {Logger};

\node (server) [rectangle, draw, thick, rounded corners=3pt, text width = 100pt, right = 2cm of internet, fill=blue!50] {Server};


\draw [thick] (app) -- node [left] {localhost:10000} (autosys);
\draw [thick] (autosys) -- node[left] {localhost:9050} (tor);

\draw [thick, color=blue]  (tor) -> (internet) {};
\draw [thick, color=blue]  (internet) -> (server) {};
\draw [thick, color=red]   (autosys) -| (internet) {};
\draw [thick, color=red]   (internet) -- node [right]{Tracker packet} (analysis);
\end{tikzpicture}
\caption {Autosys proxy structure.}
\label{fig:autosys_diagram}
\end{figure}

\subsection{Analyzer plug-in}
Data sniffed by the autosys plug-in need to be stored for later analysis, 
to do so we need another plug-in.
This plug-in is called the "Analyzer plug-in", "Analysis server" or simpy "Logger"
which will collect all the dead-drops generated by the sniffers.
The plug-in is implemented as a simple event driven UDP server, every
packet received from the sniffer plug-in is so saved in a single place.
\begin{figure}[H]
\begin{lstlisting}[language=bash,frame=single]
host_type;hostname;timestamp
\end{lstlisting}
\caption{Analyzer log Structure}
\label{fig:analyzer_pack_struct}
\end{figure}
Referring to the figure \ref{fig:analyzer_pack_struct} we have the following
information in the log:
\begin{itemize}
\item Host Type \hfill \\
\emph{c} to signal that the traced packet referred to a client\\
\emph{s} to signal that the traced packet referred to a server
\item Hostname \hfill \\
The name of the tracked machine
\item Timestamp \hfill \\
The time-stamp of the packet (when it flows trough the autosys plug-in).
\end{itemize}

In conclusion this raw data will be passed to the phase two to be analyzed by the scripts.\\
\begin{figure}[H]
\centering
\begin{tikzpicture}[scale=2]
\node (internet) [cloud, draw, thick, fill=red!20!blue!20!green!20] {Net};
\node (analysis) [rectangle, draw, thick, rounded corners=3pt, text width = 100pt, below= 1cm of internet, fill=yellow!50] {Logger};

\pic(disc2) [
	draw,
	below = 1.7cm of analysis,
	scale = 0.5
] {disc};
\pic(disc1) [
	draw,
	below = 1.1cm of analysis,
	scale = 0.5
] {disc};
\pic(disc0) [
	draw,
	below = 0.5cm of analysis,
	scale = 0.5
] {disc};


\draw [thick, color=red]   (internet) -- node [right]{Tracker packet} (analysis);
\draw [thick, color=red]   (analysis) -- (0,-1.37);
\end{tikzpicture}
\caption {Logger logic structure.}
\label{fig:autosys_diagram}
\end{figure}

As a subnote this plug-in was implemented as an UDP server for the sake of
simplicity and, in opposite to TCP, the low impact on the network.\\
This can lend us to a little packet-loss but won't impact too much on the time analysis
as we will see later on this study.

\subsection{Simpletcp plug-in}
The last plug-in that we implemented emulate a simple HTTP client/server communication.
% TODO : scrivi + socks
